% Generated draft: requires mathematical and source review.
\documentclass[11pt]{article}
\usepackage[margin=25mm]{geometry}
\usepackage{fontspec}
\setmainfont{Latin Modern Roman}
\usepackage{amsmath,amssymb,mathtools,unicode-math}
\setmathfont{Latin Modern Math}
\usepackage{enumitem,needspace,xurl,hyperref}
\setlength{\parindent}{0pt}
\setlength{\parskip}{5pt}
\setlength{\emergencystretch}{3em}
\newcommand{\N}{\mathbb{N}}
\newcommand{\Z}{\mathbb{Z}}
\newcommand{\Q}{\mathbb{Q}}
\newcommand{\R}{\mathbb{R}}
\newcommand{\C}{\mathbb{C}}
\newcommand{\F}{\mathbb{F}}
\newcommand{\A}{\mathbb{A}}
\newcommand{\PP}{\mathbb{P}}
\newcommand{\E}{\mathbb{E}}
\newcommand{\eps}{\varepsilon}
\newcommand{\dd}{\,\mathrm{d}}
\begin{document}
\section*{UnsolvedMath ID: 20003280}
\textbf{Source problem code:} catalogue entry
\subsection{Definitions and mathematical statement}
Let $S$ be a nonempty subset of the places of $\Q$. A natural precise meaning of the source's partial adele ring is
\[
 \mathbb A_S=\prod_{v\in S}'\Q_v,
\]
with restricted product taken relative to $\Z_p$ at the nonarchimedean places; $\Q_\infty=\R$ when $\infty\in S$. Put $G_S=\mathrm{SL}_2(\mathbb A_S)$ and
\[
 D_S=\left\{\begin{pmatrix}a&0\\0&a^{-1}\end{pmatrix}:a\in\mathbb A_S^\times\right\}.
\]
For a lattice $\Gamma<G_S$, E. Lindenstrauss's Question 13 asks for a classification of nonzero finite ergodic $D_S$-invariant measures on $G_S/\Gamma$, equivalently invariant ergodic probability measures after normalization. The source puts ``irreducible'' in quotation marks but does not define it, specify $S$, or specify whether the full diagonal or an identity-component subgroup is intended. For a finite product, a common irreducibility condition is dense projection to every proper nonempty subproduct; its intended extension to arbitrary $S$ needs clarification.

The supplied aid reports the following finite-place partial statements. Suppose $S$ is finite, contains $\infty$, has $|S|\ge2$, and $\Gamma$ is $S$-arithmetic, commensurable with $\mathrm{SL}_2(\mathcal O_S)$, where
\[
 \mathcal O_S=\Z[p^{-1}:p\in S\setminus\{\infty\}].
\]
It attributes to the Einsiedler--Lindenstrauss torus theorem the dichotomy that every full-$D_S$ ergodic probability $\mu$ is Haar or satisfies
\[
 h_\mu(a)=0\quad\text{for every }a\in D_S.
\]
This leaves the zero-entropy measures unclassified and is not a theorem for all subsets of places and all lattices in the original question.

For a finite product $G_S$, if a coordinate projection $\Lambda_v=\pi_v(\Gamma)$ is discrete, the aid reports that it is a lattice and that the coordinate map
\[
 G_S/\Gamma\longrightarrow \mathrm{SL}_2(\Q_v)/\Lambda_v,
 \qquad g\Gamma\longmapsto\pi_v(g)\Lambda_v
\]
sends a full-diagonal ergodic probability to a diagonal ergodic probability on the lower-place quotient. Dense coordinate projections rule out this particular discrete-projection descent, not every possible measurable factor.

For a closed diagonal orbit $D_Sg\Gamma$, its invariant orbit measure is finite exactly when the stabilizer $D_S\cap g\Gamma g^{-1}$ is a lattice in $D_S$. A cocompact such stabilizer gives a compact orbit and a finite ergodic orbit measure. The aid illustrates non-Haar measures in the one-real-place case $\mathrm{SL}_2(\R)/\mathrm{SL}_2(\Z)$ using
\[
 \gamma=\begin{pmatrix}2&1\\1&1\end{pmatrix},\qquad
 \operatorname{spec}(\gamma)=\left\{\frac{3+\sqrt5}{2},\frac{3-\sqrt5}{2}\right\}.
\]
Diagonalizing this hyperbolic element gives a periodic diagonal orbit and its normalized orbit measure. Thus Haar uniqueness already fails in this reading; the question asks for classification rather than blanket uniqueness. The appended section 3.6 sentence about algebraic actions of ``large'' groups is a transition, not an added hypothesis. These partial conclusions are supplied machine-generated research claims requiring independent verification.
\subsection{Short English statement}
Classify finite ergodic diagonal-invariant measures on partial-adelic $\mathrm{SL}_2$ quotients. The supplied finite $S$ arithmetic result separates Haar measure from zero-entropy possibilities; coordinate descent and periodic orbits explain why the choice of places and irreducibility matter.
\subsection{Sources}
\begin{itemize}
\item E. Lindenstrauss, Question 13, American Institute of Mathematics, \emph{Emerging applications of measure rigidity}, \url{https://aimath.org/WWN/measrigid/measrigid.pdf}.
\item The attached research summary attributes the finite-$S$ entropy dichotomy to the Einsiedler--Lindenstrauss torus theorem, without supplying its full bibliographic reference.
\item Local UnsolvedMath registry, permanent ID $20003280$, source code \texttt{catalogue entry}, canonical location \texttt{aim-other-notes.json}, \texttt{notes[86]}, dated 14 August 2026. Research attempt 1: partial result, proved rigor, low novelty confidence; machine-generated and not peer reviewed.
\end{itemize}
\end{document}
